\chapter{Risque de liquidité et de financement~; trésorerie bancaire}\label{ch:rc:liquidity-and-funding-risk-bank-treasury}\chaptermark{Liquidité, financement et trésorerie}

Le 8 mars 2023, la société mère de Silicon Valley Bank annonce avoir cédé 21~milliards de dollars de
titres avec une perte après impôt de 1.8~milliard et qu'elle lèverait 2.25~milliards de capital. Le
lendemain, les déposants retirent plus de 40~milliards de dollars, et la direction de la banque anticipe
plus de 100~milliards de retraits supplémentaires le 10 mars, soit environ 85~\% de ses dépôts. La banque fait défaut ce jour-là.
Elle avait investi ses dépôts, dont 94~\% n'étaient pas garantis, dans des titres hypothécaires d'agence à longue échéance
détenus jusqu'à l'échéance, avec une duration de 6.2 ans~; lorsque les taux ont monté, leur valeur a baissé et ses
déposants, un réseau étroitement lié d'entreprises technologiques et de leurs investisseurs, ont retiré leurs fonds. Rien dans
cette histoire ne relève d'un desk de trading~: il s'agit du risque de marché du bilan d'une banque et du
risque de liquidité de son financement, que gère sa trésorerie. Ce chapitre mesure les deux~: les
ratios de liquidité réglementaires, les prix de cession interne qui facturent aux métiers le financement, le
\omterm{def:rc:liquidity-and-funding-risk-bank-treasury:alm}{risque de taux du portefeuille bancaire} et le comportement des dépôts.

\section{Ratios de liquidité}

\begin{definition}[Ratio de liquidité à court terme, actifs liquides de haute qualité]\label{def:rc:liquidity-and-funding-risk-bank-treasury:lcr}
Le \emph{ratio de liquidité à court terme}\index{ratio de liquidité à court terme} (liquidity coverage ratio, LCR) est le stock de
\emph{actifs liquides de haute qualité}\index{actifs liquides de haute qualité} (high-quality liquid assets, HQLA) rapportés aux sorties nettes de trésorerie
d'un scénario de stress de 30 jours~; il doit être d'au moins 100~\%. Les HQLA sont la trésorerie, les réserves de banque centrale et
les titres de niveau~1 en totalité, les actifs de niveau~2A (tels que les titres hypothécaires d'agence) après une décote de 15~\%,
et les actifs de niveau~2 sont plafonnés à 40~\% du stock~; les sorties appliquent des taux d'écoulement à chaque source de financement,
de 5~\% pour les dépôts de particuliers stables et garantis à 40~\% pour les dépôts d'entreprises non opérationnels.
\end{definition}

\begin{definition}[Ratio de financement stable net]\label{def:rc:liquidity-and-funding-risk-bank-treasury:nsfr}
Le \emph{ratio de financement stable net}\index{ratio de financement stable net} (net stable funding ratio, NSFR) est le financement stable disponible
rapporté au financement stable requis sur un horizon d'un an, pondéré par le degré de stabilité de chaque passif et
par le degré d'illiquidité de chaque actif~; il doit être d'au moins 100~\%.
\end{definition}

\begin{example}[Une banque stylisée]\label{ex:rc:liquidity-and-funding-risk-bank-treasury:bank}
Une banque illustrative (en milliards de dollars) détient 15 de trésorerie, 25 de bons du Trésor à cinq ans, 90 de titres hypothécaires
d'agence (un échéancier à mensualités constantes sur douze ans, duration 6.3 à l'achat) et 70 de prêts à taux variable,
financés par 173 de dépôts, 11 de dette à terme et 16 de fonds propres. Un an après que les taux sont passés de 1.0~\%
à 4.5~\%, ses HQLA valorisés au prix de marché s'élèvent à 61.0~: 36.6 de niveau~1 (trésorerie et bons du Trésor) et 24.4 de niveau~2A,
les MBS après décote et plafonnement. Avec 6~\% de ses dépôts garantis et stables (écoulement 5~\%) et
94~\% non garantis (40~\%), les sorties à 30 jours s'élèvent à 65.6 et le LCR vaut 93~\%, en dessous du minimum.
\end{example}

\omcode{code/firm/alm/firm_alm.py}{122}{131}{Le LCR~: actifs de niveau~1 et de niveau~2A avec la décote et le plafond, face aux sorties pondérées par les taux d'écoulement.}\label{lst:rc:liquidity-and-funding-risk-bank-treasury:lcr}

\section{Taux de cession interne}

\begin{definition}[Taux de cession interne]\label{def:rc:liquidity-and-funding-risk-bank-treasury:ftp}
Le \emph{taux de cession interne}\index{taux de cession interne} (funds-transfer pricing, FTP) facture à chaque métier qui utilise
des financements, et crédite chaque métier qui en lève, un taux interne~: une courbe de taux de base pour l'échéance de refixation
plus une prime de liquidité pour la durée du financement, de sorte que la marge d'un desk soit mesurée
après le coût des fonds qu'il utilise, la trésorerie supportant (et couvrant) le risque de taux et
le décalage de liquidité au niveau central.
\end{definition}

\begin{example}[Une courbe de taux de cession interne]\label{ex:rc:liquidity-and-funding-risk-bank-treasury:ftp}
Avec un taux de base de 4.5~\% et une prime de liquidité illustrative de $0.15~\%\sqrt T$, un prêt de trois mois est
facturé à 4.58~\%, un prêt de cinq ans à 4.84~\% et un prêt de dix ans à 4.97~\%~; un dépôt censé rester cinq ans
est crédité à 4.84~\%. Un métier qui finance des actifs longs par des dépôts courts voit la prime de terme dans sa
facturation, au lieu d'apparaître dans les pertes de la trésorerie.
\end{example}

\section{Gestion actif--passif et risque de taux du portefeuille bancaire}

\begin{definition}[Gestion actif--passif, IRRBB]\label{def:rc:liquidity-and-funding-risk-bank-treasury:alm}
La \emph{gestion actif--passif}\index{gestion actif--passif} (asset--liability management, ALM) gère le décalage
entre les actifs et les passifs d'une banque en taux, en échéances et en devises. Son principal risque de marché est
le \emph{risque de taux du portefeuille bancaire}\index{risque de taux du portefeuille bancaire}
(interest-rate risk in the banking book, IRRBB), mesuré en valeur et en résultat.
\end{definition}

\begin{definition}[Valeur économique des fonds propres, sensibilité de la marge nette d'intérêt]\label{def:rc:liquidity-and-funding-risk-bank-treasury:eve}
La \emph{valeur économique des fonds propres}\index{valeur économique des fonds propres} (economic value of equity, EVE) est la valeur actuelle des
flux d'actifs du \omterm{def:rc:regulatory-capital-for-trading-books:books}{portefeuille bancaire}, diminuée de celle des flux de passifs (dépôts placés selon des
modèles comportementaux). La \emph{sensibilité de la marge nette d'intérêt}\index{sensibilité de la marge nette d'intérêt} (net interest income sensitivity, NII) est
la variation, sur l'année suivante, des produits d'intérêts diminués des charges d'intérêts sous un scénario de taux, à
bilan constant.
\end{definition}

La norme prudentielle prescrit six chocs instantanés par devise~: parallèles à la hausse et à la baisse,
taux courts à la hausse et à la baisse (maximal à l'extrémité courte, décroissant comme $e^{-t/4}$), un steepener et un
flattener, calibrés pour le dollar américain à 200, 300 et 150 points de base pour les mouvements parallèles, courts et longs
(voir la \cref{fig:rc:liquidity-and-funding-risk-bank-treasury:shocks}~; le graphique de la \cref{fig:rc:liquidity-and-funding-risk-bank-treasury:eve} indique les mouvements à la hausse $+$ et à la baisse $-$). Une banque dont la perte EVE maximale
dépasse 15~\% de ses fonds propres de catégorie~1 est considérée comme aberrante.

\begin{dated}{2026-09}{Les dimensions des chocs}\label{dat:rc:liquidity-and-funding-risk-bank-treasury:shocks}
En juillet 2024, le Comité de Bâle a recalibré les chocs à partir de données allant jusqu'à décembre 2023, avec des facteurs de choc
locaux, un 99.9\up{e}~centile et un arrondi à 25 points de base, pour une mise en œuvre au 1 janvier 2026. Pour le
dollar américain, les chocs parallèles et courts restent à 200 et 300 points de base et le choc long passe de 150
à 225. Les figures du chapitre utilisent les dimensions de 2016, en vigueur en 2022.
\end{dated}

\begin{omfigure}
\begin{tikzpicture}
\begin{axis}[width=11.5cm,height=5.4cm,
  xlabel={échéance (années)},ylabel={choc (pb)},
  tick label style={font=\small},label style={font=\small},
  xmin=0,xmax=20,ymin=-320,ymax=320,grid=major,grid style={gray!25},
  legend columns=3,legend cell align=left,legend style={font=\footnotesize,at={(0.5,-0.25)},anchor=north,draw=none,fill=none,/tikz/every even column/.append style={column sep=8pt}},
  table/col sep=comma]
\addplot[omDef,very thick] table[x=t,y=pu]{figdata/rates-credit-risk/24-liquidity-and-funding-risk-bank-treasury/shocks.csv};
\addplot[omDef,very thick,dashed] table[x=t,y=pd]{figdata/rates-credit-risk/24-liquidity-and-funding-risk-bank-treasury/shocks.csv};
\addplot[omThm,thick] table[x=t,y=su]{figdata/rates-credit-risk/24-liquidity-and-funding-risk-bank-treasury/shocks.csv};
\addplot[omThm,thick,dashed] table[x=t,y=sd]{figdata/rates-credit-risk/24-liquidity-and-funding-risk-bank-treasury/shocks.csv};
\addplot[omMeth,thick,mark=*,mark size=1pt] table[x=t,y=st]{figdata/rates-credit-risk/24-liquidity-and-funding-risk-bank-treasury/shocks.csv};
\addplot[omExo,thick,mark=square*,mark size=1pt] table[x=t,y=fl]{figdata/rates-credit-risk/24-liquidity-and-funding-risk-bank-treasury/shocks.csv};
\legend{choc parallèle à la hausse,choc parallèle à la baisse,choc court à la hausse,choc court à la baisse,steepener,flattener}
\end{axis}
\end{tikzpicture}

\omcaption{Les six chocs de taux prudentiels pour le dollar américain, par échéance~: parallèles
de $\pm200$ points de base, taux court de $\pm300$ décroissant avec l'échéance, et les deux rotations construites à partir des
chocs court et long. Données~: norme IRRBB du Comité de Bâle (2016)~; le
tutoriel du chapitre.}\label{fig:rc:liquidity-and-funding-risk-bank-treasury:shocks}
\end{omfigure}

\omcode{code/firm/alm/firm_alm.py}{19}{37}{Les six scénarios de chocs de taux d'intérêt prudentiels.}\label{lst:rc:liquidity-and-funding-risk-bank-treasury:shocks}

\begin{example}[La banque stylisée avant la hausse]\label{ex:rc:liquidity-and-funding-risk-bank-treasury:eve}
À un taux de 1~\%, l'EVE de la banque est de 22.9. Le choc parallèle à la hausse la réduit de 6.95, soit 43~\% de ses fonds propres Tier~1 de 16, près de trois fois le seuil de valeur atypique~;
le steepener la réduit de 4.60 (6.79 avec les montants de 2024)~;
le choc parallèle à la baisse l'augmente de 8.64 (voir la \cref{fig:rc:liquidity-and-funding-risk-bank-treasury:eve}). Le risque était
mesurable un an avant de se réaliser.
\end{example}

\begin{omfigure}
\begin{tikzpicture}
\begin{axis}[width=11.5cm,height=5.2cm,ybar,bar width=16pt,
  symbolic x coords={parallel+,parallel-,short+,short-,steepener,flattener},xtick=data,xticklabels={parallèle+,parallèle-,court+,court-,steepener,flattener},
  x tick label style={font=\footnotesize},
  ylabel={variation de l'EVE (milliards de dollars)},
  tick label style={font=\small},label style={font=\small},
  ymin=-9.5,ymax=11,grid=major,grid style={gray!25},
  nodes near coords,every node near coord/.append style={font=\footnotesize,/pgf/number format/fixed,/pgf/number format/precision=2},
  table/col sep=comma]
\addplot[fill=omDef!60,draw=omDef] table[x=scenario,y=delta]{figdata/rates-credit-risk/24-liquidity-and-funding-risk-bank-treasury/eve.csv};
\draw[omThm,dashed,thick] (axis cs:{parallel+},-2.4) -- (axis cs:{flattener},-2.4);
\end{axis}
\end{tikzpicture}

\omcaption{Variation de la \omterm{def:rc:liquidity-and-funding-risk-bank-treasury:eve}{valeur économique des fonds propres} de la banque stylisée selon les six
chocs prudentiels, à un taux de 1~\%. La ligne en pointillés figure le seuil de valeur atypique, soit une perte de 15~\% des fonds propres Tier~1 (2.4)~; les
scénarios de hausse parallèle et de steepener le franchissent. Données~: tutoriel du chapitre.}\label{fig:rc:liquidity-and-funding-risk-bank-treasury:eve}
\end{omfigure}

\section{Dépôts~: le modèle comportemental}

\begin{definition}[Dépôt sans échéance, bêta des dépôts]\label{def:rc:liquidity-and-funding-risk-bank-treasury:nmd}
Un \emph{dépôt sans échéance}\index{dépôt sans échéance} peut être retiré à tout moment et rapporte un taux
que fixe la banque~; un modèle comportemental le scinde en une part stable, censée rester et répartie sur
plusieurs années, et une part volatile. Son \emph{bêta des dépôts}\index{bêta des dépôts} mesure la part
d'une variation des taux de marché répercutée sur le taux du dépôt.
\end{definition}

Les modèles de dépôts transforment un passif juridiquement exigible à vue en un passif synthétique long~: ils justifient
le placement des dépôts dans des actifs longs. Si le bêta est faible, la hausse des taux élargit la marge et fait croître le NII~;
si les déposants exigent un bêta élevé, ou s'en vont, la marge s'inverse.

\begin{example}[Deux bêtas]\label{ex:rc:liquidity-and-funding-risk-bank-treasury:beta}
À un taux de 1~\%, la banque stylisée dégage un NII de 4.27 par an. Avec un \omterm{def:rc:liquidity-and-funding-risk-bank-treasury:nmd}{bêta des dépôts} de 0.35, une hausse de 200 points
de base ajoute 0.49 par an~; avec un bêta de 0.80, la même hausse coûte 1.07
(voir la \cref{fig:rc:liquidity-and-funding-risk-bank-treasury:nii}). Les titres à taux fixe rapportent autant
dans les deux cas~; la différence tient à ce que prennent les déposants.
\end{example}

\begin{omfigure}
\begin{tikzpicture}
\begin{axis}[width=11.5cm,height=5.2cm,
  xlabel={hausse parallèle des taux (pb)},ylabel={variation du NII (milliards de dollars par an)},
  tick label style={font=\small},label style={font=\small},
  xmin=0,xmax=400,ymin=-2.4,ymax=1.2,grid=major,grid style={gray!25},
  legend columns=2,legend style={font=\footnotesize,at={(0.5,-0.25)},anchor=north,draw=none,fill=none,/tikz/every even column/.append style={column sep=8pt}},
  table/col sep=comma]
\addplot[omMeth,very thick] table[x=bp,y=beta35]{figdata/rates-credit-risk/24-liquidity-and-funding-risk-bank-treasury/nii.csv};
\addplot[omThm,very thick,dashed] table[x=bp,y=beta80]{figdata/rates-credit-risk/24-liquidity-and-funding-risk-bank-treasury/nii.csv};
\legend{bêta des dépôts 0.35,bêta des dépôts 0.80}
\end{axis}
\end{tikzpicture}

\omcaption{Sensibilité à un an de la marge nette d'intérêt (NII) de la banque stylisée à une hausse parallèle des taux,
pour deux \omterm{def:rc:liquidity-and-funding-risk-bank-treasury:nmd}{bêtas des dépôts}. Le signe de l'exposition des revenus de la banque dépend du comportement des déposants.
Données~: tutoriel du chapitre.}\label{fig:rc:liquidity-and-funding-risk-bank-treasury:nii}
\end{omfigure}

\section{Les faillites de banques régionales en 2023}

\begin{definition}[Portefeuille détenu jusqu'à l'échéance]\label{def:rc:liquidity-and-funding-risk-bank-treasury:htm}
Un \emph{portefeuille détenu jusqu'à l'échéance}\index{portefeuille détenu jusqu'à l'échéance} regroupe des titres que la banque entend
et peut conserver jusqu'à leur échéance, comptabilisés au coût amorti~: leurs pertes de marché
n'apparaissent ni dans le résultat ni dans les fonds propres réglementaires. En vendre une partie peut obliger
à réévaluer l'ensemble du portefeuille au prix de marché.
\end{definition}

\begin{definition}[Panique bancaire]\label{def:rc:liquidity-and-funding-risk-bank-treasury:run}
Une \emph{panique bancaire}\index{panique bancaire} (bank run) est un retrait de dépôts plus rapide que ne peut honorer la banque à partir de ses
actifs liquides, porté par la crainte des déposants que d'autres retirent leurs fonds avant eux~: individuellement
rationnelle, collectivement autoréalisatrice.
\end{definition}

\begin{example}[La banque stylisée un an plus tard]\label{ex:rc:liquidity-and-funding-risk-bank-treasury:run}
Après la hausse de 350 points de base, le portefeuille de MBS détenu jusqu'à l'échéance a perdu 19.2~\% de sa valeur, soit 17.2
milliards de dollars, 108~\% des fonds propres de la banque~; les bons du Trésor ont perdu 3.9. Le jour même, la banque peut servir les retraits avec ses
liquidités et vendre ses bons du Trésor au prix de marché~: 36.1, soit 20.9~\% de ses dépôts
(voir la \cref{fig:rc:liquidity-and-funding-risk-bank-treasury:liquidity}). Au-delà, il faut mettre les MBS en garantie ou
les vendre, ce qui réalise la perte.
\end{example}

\begin{omfigure}
\begin{tikzpicture}
\begin{axis}[width=11.5cm,height=5.2cm,xbar stacked,bar width=18pt,
  ytick=\empty,xlabel={milliards de dollars},
  tick label style={font=\small},label style={font=\small},
  xmin=0,xmax=190,ymin=0.4,ymax=1.9,
  legend columns=3,legend style={font=\footnotesize,at={(0.5,-0.3)},anchor=north,draw=none,fill=none,/tikz/every even column/.append style={column sep=8pt}}]
\addplot[fill=omMeth!60,draw=omMeth] coordinates {(15,1)};
\addplot[fill=omDef!60,draw=omDef] coordinates {(21.1,1)};
\addplot[fill=omExo!30,draw=omExo] coordinates {(72.8,1)};
\draw[omThm,very thick] (axis cs:173,0.5) -- (axis cs:173,1.5) node[above,font=\footnotesize] {dépôts 173};
\draw[omThm,dashed,thick] (axis cs:36.1,0.5) -- (axis cs:36.1,1.5) node[above,font=\footnotesize] {même jour 36.1};
\legend{liquidités,bons du Trésor au prix de marché,MBS au prix de marché (pas le même jour)}
\end{axis}
\end{tikzpicture}

\omcaption{Ressources liquides de la banque stylisée un an après la hausse, à valeur de marché, face à ses
dépôts. Seules les liquidités et les bons du Trésor peuvent se transformer en paiements le jour d'une \omterm{def:rc:liquidity-and-funding-risk-bank-treasury:run}{panique bancaire}~: 20.9~\% des
dépôts. Données~: tutoriel du chapitre.}\label{fig:rc:liquidity-and-funding-risk-bank-treasury:liquidity}
\end{omfigure}

\begin{dated}{2026-09}{Mars 2023}\label{dat:rc:liquidity-and-funding-risk-bank-treasury:svb}
Selon la revue de la Réserve fédérale d'avril 2023, Silicon Valley Bank a perdu plus de 40 milliards
de dollars de dépôts le 9 mars 2023 et en attendait encore plus de 100 milliards le 10 mars, soit environ 85~\% de
sa base de dépôts~; 94~\% de ses dépôts n'étaient pas garantis à la fin de 2022, et son \omterm{def:rc:liquidity-and-funding-risk-bank-treasury:htm}{portefeuille détenu jusqu'à l'échéance}
avait une duration de 6.2~ans. Le 12 mars, la Réserve fédérale a créé le Bank Term Funding
Program, qui prête jusqu'à un an contre des titres éligibles évalués au pair, et Silicon Valley Bank et
Signature Bank ont fait l'objet d'une résolution, tous leurs déposants étant protégés. À la fin de 2022, Silicon Valley Bank avait
déclaré plus de 15 milliards de dollars de pertes latentes sur ses titres détenus jusqu'à l'échéance, soit 89~\% de ses
fonds propres de base de catégorie~1 (US Government Accountability Office). First Republic Bank, avec 229.1 milliards
de dollars d'actifs, a été fermée le 1 mai 2023 et ses dépôts ainsi que l'essentiel de ses actifs ont été cédés à JPMorgan Chase~; la
FDIC a estimé le coût pour son fonds d'assurance à environ 13 milliards de dollars.
\end{dated}

\section{Tutoriel~: un bilan soumis à des chocs}\label{tut:rc:liquidity-and-funding-risk-bank-treasury:tutorial}

\begin{tutorial}
\textbf{Objectif.} Calculer les sensibilités de la banque stylisée en \omterm{def:rc:liquidity-and-funding-risk-bank-treasury:eve}{valeur économique des fonds propres} (EVE) et en marge nette d'intérêt (NII) sous les chocs prudentiels,
son \omterm{def:rc:liquidity-and-funding-risk-bank-treasury:lcr}{ratio de liquidité à court terme} (LCR), et la sortie de fonds que ses actifs liquides peuvent couvrir. \textbf{État final~:} les résultats des
\cref{ex:rc:liquidity-and-funding-risk-bank-treasury:bank,ex:rc:liquidity-and-funding-risk-bank-treasury:eve,ex:rc:liquidity-and-funding-risk-bank-treasury:beta,ex:rc:liquidity-and-funding-risk-bank-treasury:run}
et les quatre graphiques.
\begin{tutsteps}
\item \textbf{Bilan}~: \texttt{positions()}, \texttt{deposits(beta)}, \texttt{bank(rate)}.
\item \textbf{Valeur et résultat}~: \texttt{eve\_table()}, \texttt{nii\_table()}.
\item \textbf{Liquidité}~: \texttt{lcr\_now()}, \texttt{liquid\_same\_day()}, \texttt{htm\_loss()}.
\item \textbf{Graphiques}~: \texttt{fig\_rc\_alm.py}.
\end{tutsteps}
\textbf{Pour aller plus loin.} Couvrir les MBS par un swap payeur de taux fixe de même duration et relancer les
chocs sur l'EVE~; diviser par deux la part des dépôts stables et constater que la variation de l'EVE change de signe.
\end{tutorial}

\section{Construction~: le moteur ALM}\label{bld:rc:liquidity-and-funding-risk-bank-treasury:alm}

\begin{build}
\bfield{Rôle} Le risque du \omterm{def:rc:regulatory-capital-for-trading-books:books}{portefeuille bancaire} de l'entreprise~: EVE et NII sous des scénarios prudentiels et internes,
ratios de liquidité, \omterm{def:rc:liquidity-and-funding-risk-bank-treasury:ftp}{taux de cession interne} et capacité à absorber une \omterm{def:rc:liquidity-and-funding-risk-bank-treasury:run}{panique bancaire}~; le volet bilan du moteur de risque.
\bfield{Interface} \texttt{shock\_curve(kind)}, \texttt{SCENARIOS}~; \texttt{Position},
\texttt{Deposits}, \texttt{BalanceSheet} avec \texttt{pv}, \texttt{delta\_eve}, \texttt{nii}~;
\texttt{market\_value}~; \texttt{lcr}~; \texttt{ftp\_rate}~; \texttt{run\_capacity}~;
\texttt{bond\_price\_change}, \texttt{duration}~; \texttt{eve\_ratio}.
\bfield{Règles} Courbe de base plate, avec décalages selon les scénarios~; flux annuels~; dépôts stables répartis en
montants annuels égaux~; LCR avec le cap du niveau~2 et la décote du niveau~2A~; test des valeurs aberrantes à 15~\% des fonds propres de catégorie~1.
\bfield{Tests d'acceptation} \texttt{code/firm/alm/tests/}~: formes des chocs~; variation du prix obligataire rapportée
à la duration~; un portefeuille adossé n'a aucune sensibilité à l'EVE~; le cap du niveau~2~; l'identité du \omterm{def:rc:liquidity-and-funding-risk-bank-treasury:ftp}{taux de cession interne}~; la capacité à absorber une \omterm{def:rc:liquidity-and-funding-risk-bank-treasury:run}{panique bancaire}.
\bfield{Pour aller plus loin} Une courbe de rendement complète avec base et optionnalité dans les prêts et les dépôts~; facteurs NSFR~;
bilans dynamiques~; liquidité intrajournalière.
\end{build}

\begin{omsources}
\item Board of Governors of the Federal Reserve System, \emph{Review of the Federal Reserve's
Supervision and Regulation of Silicon Valley Bank}, avril 2023.
\item Communiqué de presse de la Réserve fédérale du 12 mars 2023 sur le Bank Term Funding Program.
\item Basel Committee on Banking Supervision, \emph{The Liquidity Coverage Ratio and liquidity risk
monitoring tools} (2013); \emph{The net stable funding ratio} (2014); \emph{Interest rate risk in the
banking book} (2016).
\end{omsources}

\section{Exercices}

\begin{exercise}[$\star$]\label{exo:rc:liquidity-and-funding-risk-bank-treasury:1}
Une banque détient 30 d'actifs de niveau~1 et 40 d'actifs de niveau~2A avant décotes. Quel est son HQLA~?
\end{exercise}

\begin{exercise}[$\star$]\label{exo:rc:liquidity-and-funding-risk-bank-treasury:2}
Estimer la baisse du prix d'un portefeuille de duration 6.2 lorsque les rendements augmentent de 350~points de base, et
expliquer pourquoi la valeur exacte s'en écarte.
\end{exercise}

\begin{exercise}[$\star$]\label{exo:rc:liquidity-and-funding-risk-bank-treasury:3}
Quel est le choc sur le taux court à cinq ans, en dollars américains~?
\end{exercise}

\begin{exercise}[$\star\star$]\label{exo:rc:liquidity-and-funding-risk-bank-treasury:4}
Pourquoi la banque stylisée voit-elle son EVE augmenter lorsque les taux baissent, et que dit cela de sa position~?
\end{exercise}

\begin{exercise}[$\star\star$]\label{exo:rc:liquidity-and-funding-risk-bank-treasury:5}
Pourquoi une part élevée de dépôts non garantis constitue-t-elle un risque de liquidité, et pas seulement un problème de crédit pour les déposants~?
\end{exercise}

\begin{exercise}[$\star\star$]\label{exo:rc:liquidity-and-funding-risk-bank-treasury:6}
Comment une courbe de \omterm{def:rc:liquidity-and-funding-risk-bank-treasury:ftp}{taux de cession interne} (FTP) modifie-t-elle l'incitation d'une activité qui prête à long terme~?
\end{exercise}

\begin{exercise}[$\star\star\star$]\label{exo:rc:liquidity-and-funding-risk-bank-treasury:7}
\emph{Programmation.} Calculer la variation de NII sur un an pour une hausse de 350~points de base, avec des \omterm{def:rc:liquidity-and-funding-risk-bank-treasury:nmd}{bêtas des dépôts} de 0.35
et de 0.80.
\end{exercise}

\begin{exercise}[$\star\star\star$]\label{exo:rc:liquidity-and-funding-risk-bank-treasury:8}
\emph{Trouver la faille.} «~Nos titres sont détenus jusqu'à l'échéance et garantis par des agences publiques~;
les hausses de taux ne peuvent pas nous nuire, car nous ne les vendrons jamais.~»
\end{exercise}

\section{Problème~: mars 2023}

\begin{problem}[{Problème du week-end --- une panique bancaire visant une banque d'apparence solvable}]\label{pb:rc:liquidity-and-funding-risk-bank-treasury:1}
La banque stylisée de ce chapitre est entrée depuis un an dans la hausse des taux. Son portefeuille HTM est comptabilisé à son coût~;
ses déposants non garantis sont liés entre eux et réagissent aux mêmes informations.

\medskip\noindent\textbf{Partie I --- Le bilan.}
\begin{enumerate}
\item Donner la perte latente sur les MBS et sa part dans les fonds propres.
\item Donner la perte sur les bons du Trésor.
\item Quels seraient les fonds propres de la banque si les deux étaient réévalués au prix de marché~?
\item Pourquoi le test prudentiel sur l'EVE a-t-il déjà signalé le risque à un taux de 1~\%~?
\item Quelle couverture aurait supprimé l'essentiel de ce risque, et à quel coût pour le résultat~?
\end{enumerate}

\medskip\noindent\textbf{Partie II --- Liquidité.}
\begin{enumerate}[resume]
\item Donner le LCR et expliquer pourquoi il est inférieur à 100~\%.
\item Donner la sortie de fonds que la banque peut couvrir le jour même.
\item Comparer ce montant avec le premier jour de la \omterm{def:rc:liquidity-and-funding-risk-bank-treasury:run}{panique bancaire} subie par SVB, en part des dépôts.
\item Quelles sources de financement la banque pourrait-elle mobiliser le deuxième jour~?
\item Comment le Bank Term Funding Program a-t-il modifié les options de la banque~?
\end{enumerate}

\medskip\noindent\textbf{Partie III --- Comportement.}
\begin{enumerate}[resume]
\item Pourquoi une base de dépôts non garantis à 94~\% diffère-t-elle d'une base de particuliers de même taille~?
\item Comment un \omterm{def:rc:liquidity-and-funding-risk-bank-treasury:nmd}{bêta des dépôts} qui augmente avec les taux modifie-t-il la marge nette d'intérêt (NII) et la \omterm{def:rc:liquidity-and-funding-risk-bank-treasury:eve}{valeur économique des fonds propres} (EVE)~?
\item Que signifie pour les déposants l'annonce d'une cession de titres et d'une levée de fonds propres~?
\item Pourquoi une \omterm{def:rc:liquidity-and-funding-risk-bank-treasury:run}{panique bancaire} (bank run) peut-elle frapper une banque solvable au coût amorti~?
\item Quel ratio aurait permis de mesurer la concentration des déposants~?
\end{enumerate}

\medskip\noindent\textbf{Partie IV --- Jugement.}
\begin{enumerate}[resume]
\item Qu'aurait dû faire la trésorerie en 2021~?
\item Les pertes latentes sur le \omterm{def:rc:liquidity-and-funding-risk-bank-treasury:htm}{portefeuille détenu jusqu'à l'échéance} doivent-elles entrer dans les fonds propres~?
\item Comment faut-il valider les modèles de dépôts~?
\item Énoncer le \emph{résultat nommé}~: la perte latente sur le \omterm{def:rc:liquidity-and-funding-risk-bank-treasury:htm}{portefeuille détenu jusqu'à l'échéance} rapportée aux fonds propres et
la sortie d'un jour à partir de laquelle la banque épuise ses actifs liquides.
\item En une phrase~: qu'est-ce qui a tué la banque, le risque de marché ou le risque de liquidité~?
\end{enumerate}
\end{problem}

\section{Questions d'entretien}

\begin{interviewq}[$\star$ \iqroles{banque, risque}]\label{iq:rc:liquidity-and-funding-risk-bank-treasury:1}
À quoi protège chacun le LCR et le NSFR~? Définissez-les et précisez le risque couvert par chacun.
\end{interviewq}

\begin{interviewq}[$\star\star$ \iqroles{banque}]\label{iq:rc:liquidity-and-funding-risk-bank-treasury:2}
Qu'est-ce que le \omterm{def:rc:liquidity-and-funding-risk-bank-treasury:ftp}{taux de cession interne} (funds-transfer pricing, FTP) et pourquoi les banques en ont-elles besoin~?
\end{interviewq}

\begin{interviewq}[$\star\star$ \iqroles{risque, chercheur}]\label{iq:rc:liquidity-and-funding-risk-bank-treasury:3}
La \omterm{def:rc:liquidity-and-funding-risk-bank-treasury:eve}{valeur économique des fonds propres} (EVE) et la \omterm{def:rc:liquidity-and-funding-risk-bank-treasury:eve}{sensibilité de la marge nette d'intérêt} (NII) peuvent avoir des signes opposés. Expliquez-le et donnez un exemple.
\end{interviewq}

\begin{interviewq}[$\star\star$ \iqroles{chercheur}]\label{iq:rc:liquidity-and-funding-risk-bank-treasury:4}
Comment modéliseriez-vous les \omterm{def:rc:liquidity-and-funding-risk-bank-treasury:nmd}{dépôts sans échéance}~?
\end{interviewq}

\begin{interviewq}[$\star\star\star$ \iqroles{banque, risque}]\label{iq:rc:liquidity-and-funding-risk-bank-treasury:5}
Qu'est-ce qui a mal tourné à la Silicon Valley Bank, en termes de gestion des risques~?
\end{interviewq}

\begin{interviewq}[$\star\star\star$ \iqroles{trader, banque}]\label{iq:rc:liquidity-and-funding-risk-bank-treasury:6}
Comment couvririez-vous l'exposition de la banque à la \omterm{def:rc:liquidity-and-funding-risk-bank-treasury:eve}{valeur économique des fonds propres} (EVE), et quel effet la couverture a-t-elle sur le résultat et la comptabilité~?
\end{interviewq}
